% ============================================================================
% ภาคผนวก A — คู่มืออ้างอิงช่องป้อนข้อมูล
% ============================================================================
\appendix
\mychapter{คู่มืออ้างอิงช่องป้อนข้อมูล}

\noindent ภาคผนวกนี้รวบรวมชื่อช่องป้อนข้อมูลและฟิลด์ต่าง ๆ ที่ปรากฏในระบบ

\section{รูปแบบไฟล์และค่าที่รองรับ}

\small
\begin{longtable}{L{6cm}L{8.5cm}}
\toprule
\textbf{รายการ} & \textbf{ค่าที่รองรับ} \\
\midrule
รูปแบบไฟล์รายวิชา/ไฟล์เงื่อนไข & \texttt{XLSX} \texttt{XLS} \texttt{XLSM} \texttt{XLSB}
  \texttt{CSV} \texttt{TSV} \texttt{ODS} \texttt{FODS} \\
ระดับความพยายามในการจัดตาราง & ค่าเริ่มต้น 20{,}000 ช่วงที่ยอมรับ 1 ถึง 1{,}000{,}000 \\
\bottomrule
\end{longtable}

\section{ช่องป้อนข้อมูลในหน้าข้อมูลและการตั้งค่า}

\small
\begin{longtable}{L{5cm}L{9.5cm}}
\toprule
\textbf{ช่อง} & \textbf{คำอธิบาย} \\
\midrule
\texttt{เพิ่มไฟล์รายวิชา} & เพิ่มไฟล์ข้อมูลรายวิชา (รองรับหลายรูปแบบ) \\
\texttt{เลือกไฟล์เงื่อนไข} & เพิ่มไฟล์เงื่อนไขการสอบ \\
\texttt{วันเริ่มสอบกลางภาค} & วันที่เริ่มช่วงสอบกลางภาค \\
\texttt{วันสิ้นสุดสอบกลางภาค} & วันที่สิ้นสุดช่วงสอบกลางภาค \\
\texttt{วันเริ่มสอบปลายภาค} & วันที่เริ่มช่วงสอบปลายภาค \\
\texttt{วันสิ้นสุดสอบปลายภาค} & วันที่สิ้นสุดช่วงสอบปลายภาค \\
\texttt{วันเสาร์–อาทิตย์} & นโยบายสำหรับวันเสาร์–อาทิตย์ \\
\texttt{วันหยุด} & นโยบายสำหรับวันหยุด \\
\texttt{เพิ่มวันหยุด} & เพิ่มวันที่เป็นวันหยุด \\
\texttt{ช่วงเวลาและการค้นหา} & ขยายตัวเลือกขั้นสูง (ช่วงเวลา สอบเต็มวัน ความพยายาม) \\
\bottomrule
\end{longtable}

\needspace{7\baselineskip}
\section{ฟิลด์ในไฟล์ข้อมูลรายวิชา}

\small
\begin{longtable}{L{4.5cm}L{10cm}}
\toprule
\textbf{ฟิลด์} & \textbf{คำอธิบาย} \\
\midrule
\texttt{courseCode} & รหัสวิชา เก้าหลัก (บังคับ) \\
\texttt{courseName} & ชื่อวิชา (บังคับ) \\
\texttt{sectionNumber} & หมายเลขตอนเรียน (บังคับ) \\
\texttt{studentGroups} & กลุ่มนักศึกษา (บังคับ) \\
\texttt{enrolled} & จำนวนผู้ลงทะเบียน (ไม่บังคับ) \\
\texttt{capacity} & จำนวนที่นั่ง (ไม่บังคับ) \\
\texttt{instructor} & ชื่ออาจารย์ (ไม่บังคับ) \\
\texttt{teachingHours} & ชั่วโมงสอน (ไม่บังคับ) \\
\bottomrule
\end{longtable}

\section{ค่าที่รองรับในไฟล์เงื่อนไขการสอบ}

ไฟล์เงื่อนไขเป็นตารางที่แถวแรกเป็นหัวตารางชื่อเงื่อนไข และมีรหัสวิชาเป็นค่าในช่องใต้หัวคอลัมน์
(ดูรูปแบบของไฟล์และตัวอย่างในบทที่ 2) หัวคอลัมน์ที่ระบบรู้จักมีดังนี้:

\small
\begin{longtable}{L{4.5cm}L{10cm}}
\toprule
\textbf{หัวคอลัมน์ในไฟล์} & \textbf{ความหมาย} \\
\midrule
\texttt{ไม่มีสอบกลางภาค} & ระบบจะไม่จัดสอบกลางภาคให้วิชาที่ระบุ \\
\texttt{ไม่มีสอบปลายภาค} & ระบบจะไม่จัดสอบปลายภาคให้วิชาที่ระบุ \\
\texttt{ไม่มีสอบ} & ระบบจะไม่จัดสอบทั้งสองภาคให้วิชาที่ระบุ \\
\texttt{สอบสองช่วง} & สอบเต็มวัน \\
\texttt{สอบตรงกัน} & จัดสอบพร้อมกันกับวิชาอื่น \\
\bottomrule
\end{longtable}

